\subsection{SUBGROUPS OF G CONTAINING B}

For any subset \(X\) of \(S\) , we denote by \(W_X\) the subgroup of \(W\) generated by \(X\) (cf.~§1, no. 8) and by \(G_X\) the union \(BW_X B\) of the double cosets \(C(w)\) , \(w \in W_X\) . We have \(G_\varnothing = B\) and \(G_S = G\) .

THEOREM 3. a) For any subset \(X\) of \(S\) , the set \(G_X\) is a subgroup of \(G\) , generated by \(\bigcup_{s \in X} C(s)\) .

\begin{enumerate}
\def\labelenumi{\alph{enumi})}
\setcounter{enumi}{1}
\item
  The map \(\mathrm{X} \mapsto \mathrm{G}_{\mathrm{X}}\) is a bijection from \(\mathfrak{P}(\mathrm{S})\) to the set of subgroups of \(\mathrm{G}\) containing \(\mathrm{B}\) .
\item
  Let \((\mathrm{X}_i)_{i\in I}\) be a family of subsets of X. If \(X = \bigcap_{i\in I}X_i\) , then \(G_X = \bigcap_{i\in I}G_{X_i}\) .
\item
  Let X and Y be two subsets of S. Then \(G_{X} \subset G_{Y}\) (resp. \(G_{X} = G_{Y}\) ) if and only if \(X \subset Y\) (resp. X = Y).
\end{enumerate}

It is clear that \(G_{X} = (G_{X})^{-1}\) ; Lemma 1 of no. 1 shows that \(G_{X}.G_{X} \subset G_{X}\) ; and hence \(a)\) follows, taking into account Cor. 1 of Th. 2.

The injectivity of \(\mathrm{X} \mapsto \mathrm{G}_{\mathrm{X}}\) follows from that of \(\mathrm{X} \mapsto \mathrm{W}_{\mathrm{X}}\) (§1, no. 8, Th. 2). Conversely, let \(\mathrm{H}\) be a subgroup of \(\mathrm{G}\) containing \(\mathrm{B}\) . Let \(\mathrm{U}\) be the set of \(w \in \mathrm{W}\) such that \(C(w) \subset \mathrm{H}\) . We have \(\mathrm{H} = \mathrm{BUB}\) since \(\mathrm{H}\) is a union of double cosets. Let \(\mathrm{X} = \mathrm{U} \cap \mathrm{S}\) ; we show that \(\mathrm{H} = \mathrm{G}_{\mathrm{X}}\) . Clearly, \(\mathrm{G}_{\mathrm{X}} \subset \mathrm{H}\) . On the other hand, let \(u \in \mathrm{U}\) and let \((s_1, \ldots, s_q)\) be a reduced decomposition of \(u\) . Cor. 3 of Th. 2 implies that \(C(s_j) \subset \mathrm{H}\) , and hence that \(s_j \in \mathrm{X}\) for \(1 \leqslant j \leqslant q\) . Thus, \(u \in \mathrm{W}_{\mathrm{X}}\) , and since \(\mathrm{H}\) is the union of the \(C(u)\) for \(u \in \mathrm{U}\) , we have \(\mathrm{H} \subset \mathrm{G}_{\mathrm{X}}\) , which proves \(b\) ).

Assertions \(c)\) and \(d)\) follow from analogous properties of \(W_X\) ( \(\S 1\) , no. 8, Th. 2).

COROLLARY. The set \(S\) consists of the elements \(w \in W\) such that \(w \neq 1\) and \(B \cup C(w)\) is a subgroup of \(G\) .

The elements \(w \in W\) such that \(B \cup C(w)\) is a subgroup of G are those for which there exists \(X \subset S\) with \(W_{X} = \{1, w\}\) . Moreover, if \(w \neq 1\) , we necessarily have \(\text{Card}(X) = 1\) , i.e.~\(w \in S\) .

Remark. 1) The above corollary shows that S is determined by (G, B, N); for this reason, we sometimes allow ourselves to say that (G, B, N) is a Tits system, or that (B, N) is a Tits system in G.

PROPOSITION 1. Let X be a subset of S and N' a subgroup of N whose image in W is equal to W\_X. Then, (G\_X, B, N', X) is a Tits system.

We have \(G_{X} = BW_{X}B = BN'B\) , which shows that \(G_{X}\) is generated by \(B \cup N'\) . The verification of the axioms (T1) to (T4) is now immediate.

PROPOSITION 2. Let \(\mathrm{X}, \mathrm{Y} \subset \mathrm{S}\) and \(w \in \mathrm{W}\) . We have

\[
\mathrm{G} _ {\mathrm{X}} w \mathrm{G} _ {\mathrm{Y}} = \mathrm{BW} _ {\mathrm{X}} w \mathrm{W} _ {\mathrm{Y}} \mathrm{B}.
\]

Let \(s_1, \ldots, s_q \in X\) and \(t_1, \ldots, t_q \in Y\) . Lemma 1 shows that

\[
\mathrm{C} (s _ {1} \dots s _ {q}). \mathrm{C} (w). \mathrm{C} (t _ {1} \dots t _ {q}) \subset \mathrm{BW} _ {\mathrm{X}} w \mathrm{W} _ {\mathrm{Y}} \mathrm{B},
\]

and hence that

\[
\mathrm{G} _ {\mathrm{X}} w \mathrm{G} _ {\mathrm{Y}} \subset \mathrm{BW} _ {\mathrm{X}} w \mathrm{W} _ {\mathrm{Y}} \mathrm{B}.
\]

The opposite inclusion is obvious.

Remark. 2) Denote by \(G_X \backslash G / G_Y\) the set of subsets of \(G\) of the form \(G_X gG_Y\) , \(g \in G\) ; and define \(W_X \backslash W / W_Y\) analogously. The preceding proposition shows that the canonical bijection \(w \mapsto C(w)\) from \(W\) to \(B \backslash G / B\) defines by passage to the quotient a bijection \(W_X \backslash W / W_Y \to G_X \backslash G / G_Y\) .

PROPOSITION 3. Let \(\mathbf{X} \subset \mathbf{S}\) and \(g \in \mathbf{G}\) . The relation \(g\mathrm{B}g^{-1} \subset \mathbf{G}_{\mathbf{X}}\) implies that \(g \in \mathbf{G}_{\mathbf{X}}\) .

Let \(w \in W\) be such that \(g \in C(w)\) . Since B is a subgroup of \(G_X\) , the hypothesis \(gBg^{-1} \subset G_X\) implies that \(C(w).C(w^{-1}) \subset G_X\) , and hence that \(C(w) \subset G_X\) by Cor. 3 of Th. 2, so g belongs to \(G_X\) .

